\documentclass[11pt,a4paper]{article}
\usepackage[margin=2.2cm]{geometry}
\usepackage[T1]{fontenc}
\usepackage{lmodern}
\usepackage{amsmath,amssymb,amsthm,mathtools}
\usepackage{enumitem}
\usepackage{fancyhdr}
\usepackage[hidelinks]{hyperref}
\usepackage{microtype}

\setlength{\parindent}{0pt}
\setlength{\parskip}{0.4em}
\pagestyle{fancy}
\setlength{\headheight}{14pt}
\fancyhf{}
\lhead{Olympiad Practice Examination --- Answer Key}
\rhead{\thepage}
\renewcommand{\headrulewidth}{0.4pt}
\newcommand{\RR}{\mathbb{R}}
\newcommand{\solution}[2]{\subsection*{Problem #1 --- #2}}

\title{\vspace{-1.5cm}\bfseries Olympiad Practice Examination\\Answer Key and Worked Solutions}
\author{}
\date{}

\begin{document}
\maketitle
\vspace{-2em}

\solution{1}{Polynomials}
The answers are
\[
  \boxed{P(x)=c\,x(x-1)(x+1)\qquad(c\in\RR\setminus\{0\}).}
\]

Substituting $x=1,-2,0$ in turn into the given identity gives
$P(1)=P(-1)=P(0)=0$. Hence
\[
  P(x)=x(x-1)(x+1)Q(x)
\]
for some real polynomial $Q$. Since
$P(x+1)=(x+1)x(x+2)Q(x+1)$, substitution and cancellation of the nonzero
polynomial $x(x-1)(x+1)(x+2)$ give
\[
  Q(x+1)=Q(x).
\]
Thus $Q(n)=Q(0)$ for every nonnegative integer $n$. The polynomial
$Q(x)-Q(0)$ therefore has infinitely many roots, so $Q$ is constant.
Nonconstancy of $P$ forces this constant to be nonzero.

Conversely, for $P(x)=c x(x-1)(x+1)$,
\[
  (x-1)P(x+1)=c(x-1)x(x+1)(x+2)=(x+2)P(x),
\]
so every displayed polynomial is a solution.

\solution{2}{Inequalities}
Set $x=1/a$, $y=1/b$, $z=1/c$. Then $xyz=1$ and
\[
  \frac1{a^3(b+c)}
  =\frac{x^3yz}{y+z}
  =\frac{x^2}{y+z}.
\]
The other two terms transform similarly. By Cauchy--Schwarz (Engel form),
\[
  \sum_{\mathrm{cyc}}\frac{x^2}{y+z}
  \ge \frac{(x+y+z)^2}{2(x+y+z)}
  =\frac{x+y+z}{2}
  \ge\frac32,
\]
where the last step is AM--GM together with $xyz=1$. Equality requires
$x=y=z$ and $xyz=1$, so equality holds exactly at
\[
  \boxed{a=b=c=1}.
\]

\solution{3}{Functional Equations}
The solutions are
\[
  \boxed{f(x)=a\sin x+b\cos x\qquad(a,b\in\RR).}
\]

Let $b=f(0)$. Putting $x=0$ gives
\[
  f(y)+f(-y)=2b\cos y. \tag{1}
\]
Define $g(x)=f(x)-b\cos x$. Since $b\cos x$ itself satisfies the given
equation, $g$ also satisfies it. Moreover, (1) shows that $g$ is odd.

Applying the equation for $g$ first with $(x,y)$ and then with $(y,x)$, and
using oddness in the second equation, yields
\begin{align*}
  g(x+y)+g(x-y)&=2g(x)\cos y,\\
  g(x+y)-g(x-y)&=2g(y)\cos x.
\end{align*}
Adding and dividing by $2$, we obtain
\[
  g(x+y)=g(x)\cos y+g(y)\cos x. \tag{2}
\]
Let $a=g(\pi/2)$. Putting $y=\pi/2$ in (2) gives
$g(x+\pi/2)=a\cos x$. Replacing $x$ by $x-\pi/2$ therefore gives
$g(x)=a\sin x$, and hence $f(x)=a\sin x+b\cos x$.

Finally, the addition formulas for sine and cosine show directly that every
function of this form satisfies the original equation.

\solution{4}{Complex Numbers}
The numbers $\omega^k$, $0\le k<n$, are precisely the roots of $z^n-1$, so
\[
  z^n-1=\prod_{k=0}^{n-1}(z-\omega^k).
\]
Taking the logarithmic derivative at $z=2$ gives
\[
  \sum_{k=0}^{n-1}\frac1{2-\omega^k}
  =\frac{n2^{n-1}}{2^n-1}.
\]
Thus the requested value is
\[
  \boxed{\dfrac{n2^{n-1}}{2^n-1}}.
\]

\solution{5}{Combinatorics}
There are $\binom{2n}{n}$ unrestricted sequences containing $n$ copies of
each sign. Call a sequence \emph{bad} if some initial segment contains more
$-$ signs than $+$ signs.

In a bad sequence, take the first position at which the running total (with
$+$ counted as $1$ and $-$ as $-1$) equals $-1$. Reflect every symbol
\emph{after} that position, interchanging $+$ and $-$. The resulting sequence
has $n-1$ copies of $+$ and $n+1$ copies of $-$.

Conversely, every sequence with these new counts must reach running total
$-1$ at some first position. Reflecting the suffix after that position
reconstructs a bad sequence with $n$ copies of each sign. These operations
are inverse bijections. Consequently, the number of bad sequences is
$\binom{2n}{n+1}$, and the number of balanced sequences is
\[
  \binom{2n}{n}-\binom{2n}{n+1}
  =\binom{2n}{n}\left(1-\frac n{n+1}\right)
  =\boxed{\frac1{n+1}\binom{2n}{n}}.
\]

\solution{6}{Pigeonhole Principle}
The maximum is $\boxed{n}$.

Every positive integer has a unique representation as $2^j q$ with $q$ odd.
For each odd $q\le 2n-1$, form the chain
\[
  C_q=\{q,2q,4q,\ldots\}\cap\{1,2,\ldots,2n\}.
\]
These $n$ chains partition the universe, and within any one chain, every two
distinct elements are comparable by divisibility. If $S$ had at least $n+1$
elements, the pigeonhole principle would put two elements in the same chain.
One would then divide the other, a contradiction. Therefore $|S|\le n$.

This bound is attained by
\[
  S=\{n+1,n+2,\ldots,2n\}.
\]
Indeed, if $a,b\in S$ and $a<b$, then $2a>2n\ge b$, so $b/a$ cannot be an
integer greater than $1$. Thus no element divides a different element, and
the maximum is exactly $n$.

\solution{7}{Graph Theory}
A connected graph in which every vertex has even degree has an Euler circuit.
Traverse one such circuit and orient each edge in the direction of traversal.
Whenever the circuit enters a vertex, it next leaves that vertex, contributing
one to its indegree and one to its outdegree. At the starting vertex, the
final arrival and initial departure are paired similarly. Hence indegree and
outdegree are equal at every vertex.

\solution{8}{Game Theory}
The losing positions are precisely the multiples of $3$.

Indeed, no power of $2$ is divisible by $3$, so every move from a multiple of
$3$ leaves a position not divisible by $3$. Conversely, from $n\equiv1\pmod3$,
removing $1$ stone leaves a multiple of $3$; from $n\equiv2\pmod3$, removing
$2$ stones leaves a multiple of $3$. These moves are legal, including at
$n=1,2$, where they win immediately. Induction on the number of stones proves
the classification.

Since $2026\equiv1\pmod3$, the \textbf{first player} has a winning strategy.
Remove $1$ stone initially. After each opposing move, the position is not a
multiple of $3$: remove $1$ stone if its residue is $1$, and $2$ stones if its
residue is $2$. This always returns a multiple of $3$ to the opponent. The
pile strictly decreases, and eventually the opponent receives $3$ stones;
after any legal move, the first player removes the remaining $1$ or $2$
stones and wins.

\solution{9}{Number Theory}
The answers are $\boxed{n=1,2,4,6}$.

Modulo $2^n+n$, we have $2^n\equiv-n$. Therefore
\[
  8^n+n=(2^n)^3+n\equiv-n^3+n=-n(n-1)(n+1).
\]
Thus the desired divisibility is equivalent to
\[
  2^n+n\mid n^3-n. \tag{1}
\]

For $n\ge10$, we claim that $2^n>n^3$. At $n=10$, this is $1024>1000$. If it
holds for some $n\ge10$, then
\[
  2^{n+1}>2n^3>(n+1)^3,
\]
because $(1+1/n)^3\le(11/10)^3<2$. This proves the claim by induction.
Consequently, for $n\ge10$,
\[
  2^n+n>n^3>n^3-n>0,
\]
so (1) is impossible. Checking the remaining values gives
\[
\begin{array}{c|c|c|c}
 n&2^n+n&n^3-n&\text{Divisibility?}\\ \hline
 1&3&0&\text{yes}\\
 2&6&6&\text{yes}\\
 3&11&24&\text{no}\\
 4&20&60&\text{yes}\\
 5&37&120&\text{no}\\
 6&70&210&\text{yes}\\
 7&135&336&\text{no}\\
 8&264&504&\text{no}\\
 9&521&720&\text{no}
\end{array}
\]
which completes the proof.

\clearpage
\solution{10}{Geometry --- Analytic}
For any point $X=(x,x^2)$ on the parabola,
\[
  XF=\sqrt{x^2+\left(x^2-\frac14\right)^2}
  =x^2+\frac14,
\]
which is its distance to the line $d:y=-\frac14$.

A vertical line through $F$ meets the parabola only once, so the given line
is nonvertical. Write it as $y=mx+\frac14$. Its intersections with the
parabola have abscissae satisfying
\[
  x^2-mx-\frac14=0.
\]
The discriminant is $m^2+1>0$, and the product $uv$ of the two roots $u,v$
is $-1/4<0$. Hence $F$, whose abscissa is $0$, lies between $A$ and $B$ on
the focal chord, giving $AF+BF=AB$.

Let $O$ be the midpoint of $AB$. The circle with diameter $AB$ has center $O$
and radius $AB/2$. Using the focal-distance calculation above, the distance
from $O$ to $d$ is
\[
  \frac{y_A+y_B}{2}+\frac14
  =\frac{(y_A+\frac14)+(y_B+\frac14)}2
  =\frac{AF+BF}{2}
  =\frac{AB}{2}.
\]
The center's distance to $d$ therefore equals the radius, proving that the
circle is tangent to $d$.

\solution{11}{Geometry --- Synthetic}
All angles in this argument are directed modulo $180^\circ$. Because
$PM\perp AB$ and $PD\perp AC$, rotating both $AB$ and $AC$ through $90^\circ$
gives
\[
  \angle MPD=\angle(PM,PD)=\angle(AB,AC)=\angle BAC.
\]
Using that $A,B,C,D$ are concyclic and that $P\in BD$, $M\in CD$, we get
\[
  \angle BAC=\angle BDC=\angle(DP,DM)=\angle PDM.
\]
Thus $\angle MPD=\angle PDM$, so $MP=MD$.

Similarly, $PM\perp AB$ and $PC\perp PB$ imply
\[
  \angle MPC=\angle(PM,PC)=\angle(AB,PB)=\angle ABD.
\]
Cyclicity and collinearity then give
\[
  \angle ABD=\angle ACD=\angle(CP,CM)=\angle PCM.
\]
Hence $\angle MPC=\angle PCM$, so $MP=MC$. Combining the two length
equalities yields $MC=MD$. Since $M$ lies on line $CD$ and is equidistant from
$C$ and $D$, it must be their midpoint. \qed

\end{document}
